% ============================================================
%  resume.tex — Résumé en français
% ============================================================
\chapter*{Résumé}
\addcontentsline{toc}{chapter}{Résumé}

La classification des sous-types moléculaires du cancer
du sein à partir de l'IRM dynamique avec contraste
(DCE-IRM) est une tâche cliniquement importante~: le
sous-type détermine si une patiente peut bénéficier d'une
hormonothérapie ou nécessite une chimiothérapie.
L'entraînement d'un classifieur automatique se heurte
cependant à plusieurs obstacles concrets~: un seul hôpital
dispose rarement de plus de 800 patientes, le déséquilibre
entre sous-types atteint 10~pour~1, et les réglementations
de protection des données --- dont la loi algérienne 18-07
--- empêchent les institutions de regrouper leurs archives
d'imagerie.

Ce travail étudie si l'apprentissage fédéré peut produire
un classifieur utile dans ce contexte, sans qu'aucune
donnée patient ne quitte son institution d'origine.
L'étude est menée sur une cohorte publique de 737~patients
(Hôpital Duke, Zenodo DOI~10.5281/zenodo.7956360) avec
quatre sous-types~: Luminal~A, Luminal~B, HER2-enrichi
et Triple-Négatif.

Un pipeline d'entraînement local en trois étapes a été
construit à partir de composants éprouvés dans la
littérature~: un backbone ConvNeXt-Nano avec apprentissage
multi-instance à attention gated (MIL) traite les volumes
entiers sans masques de segmentation~; la perte LDAM,
l'optimisation ASAM, la moyenne de poids stochastique
(SWA) et le réentraînement du classifieur (cRT) corrigent
le déséquilibre de classe à chaque étape.
Ce pipeline atteint un macro~F1 de~0,342 (quatre classes)
et 0,667 (binaire) sur le pli de validation.

Pour l'expérience fédérée, trois clients hospitaliers
simulés sont construits selon une partition de Dirichlet
($\alpha = 0{,}5$).
FedAvg s'effondre vers la classe majoritaire
(macro~F1~= 0,429).
L'algorithme proposé FedSCRT --- qui combine un backbone
SCAFFOLD et un réentraînement fédéré du classifieur avec
échantillonnage équilibré --- atteint un macro~F1 de~0,629
et un AUC de~0,687, comblant 84,1~\% de l'écart de
performance fédérée et dépassant l'AUC de référence
centralisée.
Un résultat théorique montre que l'échantillonnage équilibré
rend le gradient du classifieur indépendant de la
distribution locale des classes, et une expérience contrôlée
de similarité cosinus sur trois hôpitaux artificiels le
confirme~: les gradients équilibrés s'alignent sur l'objectif
global ($\cos \geq 0,99$) quel que soit le préalable local,
tandis que les gradients naturels dérivent vers la classe
majoritaire, expliquant la robustesse de l'étape de
réentraînement fédéré du classifieur face à l'hétérogénéité
des distributions hospitalières.

\vspace{1em}
\noindent\textbf{Mots-clés~:} Apprentissage fédéré,
IRM mammaire, classification moléculaire, déséquilibre
de classes, FedSCRT, MIL à attention gated, LDAM,
SCAFFOLD.
